% ===================================================================
%  TAFEL Nr. 3 — Kindheit und Jugend.
%  Traduction 2026 d'apres texto/io, controlee sur texto/fr et
%  texto/en. Consigne : voir texto/de/10-tafel-01.tex.
%
%  Deux scenes, et l'ido subdivise beaucoup plus que le francais :
%  huit des intertitres n'ont pas de vis-a-vis francais. L'allemand
%  suit l'ido, comme partout ailleurs.
%
%  CETTE PAGE EST LE REVERS EXACT DE LA PRECEDENTE, ET C'EST LA
%  MEILLEURE CHOSE QUE L'ALLEMAND AIT A DIRE.
%
%  Le tableau 2 a montre Jahn gagnant : il a refuse « Gymnastik » en
%  1811 et impose « Turnen », qui tient encore. Le tableau 3 montre la
%  MEME campagne, menee par les memes gens, perdue de bout en bout.
%
%  L'ALLEMAND A EU DES NOMS DE MOIS A LUI, et l'alinea 1 les aurait
%  portes : Lenzing pour mars, Ostermond pour avril, Wonnemond pour
%  mai ; et a la seconde scene Brachet pour juin, HEUERT pour juillet,
%  Ernting pour aout. Charlemagne les avait fixes, les puristes du
%  XVIIIe et du XIXe siecles ont essaye deux fois de les rendre a
%  l'usage, et l'usage a garde März, April, Mai, Juni, Juli, August.
%  La colonne ecrit donc les noms latins, parce que ce sont ceux
%  qu'on dit.
%
%  ET « HEUERT » VEUT DIRE LE MOIS DES FOINS. Le tableau 3 estonien
%  avait trouve « heinakuu », le mois des foins ; le tableau 3
%  romanche « fanadur », le faneur. Trois langues sans rapport ont
%  nomme juillet par le foin, et l'allemand est la seule des trois qui
%  ait laisse tomber son mot. Le renvoi de gravuri/korekti.json qui
%  met « Juli » en gras porte donc sur le seul mot de la page ou
%  l'allemand ait perdu ce que l'estonien et le romanche ont garde.
%
%  DEUX CAMPAGNES, DEUX RESULTATS, ET LA DIFFERENCE SE MESURE :
%  Jahn nommait des choses NEUVES — un agres qui venait d'etre
%  construit n'avait pas encore de nom, et le premier nom donne
%  restait. Les puristes du calendrier voulaient reprendre un nom a
%  une chose qui en avait deja un depuis mille ans. ON PEUT DECIDER
%  DU NOM D'UNE CHOSE QUI ARRIVE ; ON NE PEUT PAS REPRENDRE LE NOM
%  D'UNE CHOSE QUI EST LA.
%
%  LA MEME PERTE, UNE SECONDE FOIS, A L'ALINEA 19 : le cousin et la
%  cousine. L'allemand disait « Vetter » et « Base », mots germaniques
%  attestes depuis le VIIIe siecle. Le francais les a remplaces au
%  XIXe — « Cousin », « Cousine » — et « Base » est mort, « Vetter »
%  survit dans les campagnes. La colonne ecrit « Vetter » et « Base »,
%  qui sont les mots du livret et se comprennent encore, mais c'est le
%  seul endroit des seize tableaux ou elle choisisse le mot d'hier :
%  ailleurs, la consigne de 2026 l'interdit. On l'ecrit ici pour que
%  personne ne prenne cette exception pour un oubli.
%
%  ENFIN UN COMPOSE QUI A GAGNE, ET IL EST DU MEME AGE QUE LE
%  FAHRRAD : l'Achterbahn de l'alinea 13, la voie-en-huit,
%  pour ce que le francais appelle montagnes russes et l'anglais
%  d'alors « switchback ». Trois langues, trois images, et aucune
%  n'est celle de l'autre — la Russie, l'aiguillage, le huit. La chose
%  est arrivee de la foire dans les annees 1880, personne ne l'avait
%  nommee, et chacun l'a nommee pour soi. C'est exactement la
%  condition de Jahn.
% ===================================================================

%%K t03-apar-1 apar
\VUsaut{3.0mm}
\VUtitre{15.0pt}{60}{TAFEL Nr. 3}
\VUsaut{1.6mm}
\VUtitre{11.4pt}{0}{Kindheit und Jugend.}
\VUsaut{3.4mm}

%%K t03-apar-2 apar
(Die Tafeln 3 und 4 sind so angelegt, dass sie in vier
\VUgras{Familienszenen} den wichtigsten Wortschatz zu \VUgras{Wetter},
\VUgras{Lebensalter}, \VUgras{Familie}, \VUgras{Kleidung} und
\VUgras{Gesundheit} bringen.)
\VUsaut{4.0mm}

%%K t03-tit-1 sub
\VUcentre{12.0pt}{0}{\textit{Erste Szene.}}
\VUsaut{3.0mm}

%%K t03-tit-2 sub
\VUpk{10.2pt}{40}{Eine Taufe im Dorf.}
\VUsaut{2.4mm}

%%K t03-tit-3 sub
\VUpk{10.2pt}{40}{Frühling.}
\VUsaut{3.0mm}

%%K t03-c1-01-1 p
1. --- Es ist \VUgras{Frühling}, im Monat \VUgras{März}, \VUgras{April} oder
\VUgras{Mai}, an einem Tag der \VUgras{Woche} --- \VUgras{Montag},
\VUgras{Dienstag} oder \VUgras{Mittwoch}. Es ist \VUgras{zehn Uhr} morgens.

%%K t03-c1-02-1 p
2. --- Das \VUgras{Wetter} ist schön, die \VUgras{Luft} klar. Die Sonne scheint.
Ein paar kleine \VUgras{Wolken} \textsuperscript{(2)} ziehen am blauen
\VUgras{Himmel} \textsuperscript{(1)} herauf.

%%K t03-c1-03-1 p
3. --- Die \VUgras{Bäume} haben noch kaum \VUgras{Blätter}. Die schlanken
\VUgras{Pappeln} \textsuperscript{(3)} und die hohe \VUgras{Platane}
\textsuperscript{(17)} tragen erst Knospen.

%%K t03-c1-04-1 p
4. --- Die \VUgras{Schwalben} \textsuperscript{(4)} sind zurück und kreisen mit
fröhlichem Geschrei um die \VUgras{Turmspitze} \textsuperscript{(7)} des
\VUgras{Glockenturms} \textsuperscript{(5)}, der die \VUgras{Dorfkirche}
\textsuperscript{(6)} krönt.

%%K t03-tit-4 sub
\VUsaut{3.4mm}
\VUpk{10.2pt}{40}{Die Kirche.}
\VUsaut{3.0mm}

%%K t03-c1-05-1 p
5. --- Diese Kirche ist sehr schlicht, mit ihrem großen \VUgras{Portal}
\textsuperscript{(13)} und ihrer großen \VUgras{Uhr} \textsuperscript{(8)}. Der kleine
\VUgras{Zeiger} \textsuperscript{(9)} steht auf der X: er zeigt die
\VUgras{Stunden} an. Der \VUgras{große} \textsuperscript{(10)} steht auf der XII
(Mittag); er zeigt die \VUgras{Minuten} an.

%%K t03-c1-06-1 p
6. --- Im Turm läuten die \VUgras{Glocken} \textsuperscript{(11)} fröhlich. Oben auf
dem Dach steht ein kleines steinernes \VUgras{Kreuz} \textsuperscript{(12)}.

%%K t03-c1-07-1 p
7. --- Die Kirche hat nicht nur ein großes \VUgras{Portal} \textsuperscript{(13)},
sondern auch eine kleine Seitentür. Durch diese Tür kommt der \VUgras{Pfarrer}
\textsuperscript{(15)} in seiner \VUgras{Soutane} aus der \VUgras{Sakristei}
\textsuperscript{(14)} und geht zum \VUgras{Pfarrhaus} \textsuperscript{(19)} hinüber.

%%K t03-c1-08-1 p
8. --- Neben ihm steht die \VUgras{Milchfrau} \textsuperscript{(24)} mit ihrem
\VUgras{Esel} \textsuperscript{(23)}. Sie kommt aus der Stadt zurück, wo sie ihre
gute Milch für die Kinder abgeliefert hat.

%%K t03-tit-5 sub
\VUsaut{4.0mm}
\VUpk{10.2pt}{40}{Der Obstgarten.}
\VUsaut{3.4mm}

%%K t03-c1-09-1 p
9. --- Meister Grauohr (der Esel) ist mit diesem Morgenlauf ganz zufrieden. Mit
Behagen atmet er die Luft ein, die nach den \VUgras{Blüten} \textsuperscript{(20)}
duftet, die die \VUgras{Obstbäume} \textsuperscript{(18)} des benachbarten
\VUgras{Obstgartens} bedecken. Die \VUgras{Pfirsichbäume} \textsuperscript{(18)} und
die \VUgras{Aprikosenbäume} \textsuperscript{(19)} stehen ganz in Rosa und Weiß und
versprechen eine gute Ernte.

%%K t03-c1-10-1 p
10. --- Unterdessen fliegen die kleinen \VUgras{Bienen} \textsuperscript{(22)} aus dem
\VUgras{Bienenstock} \textsuperscript{(21)} dorthin und sammeln summend ihren süßen
\VUgras{Honig}.

%%K t03-c1-11-1 p
11. --- Aber was ist da los? Die Dorfkinder kommen angerannt und scheuchen die
erschrockenen \VUgras{Gänse} \textsuperscript{(25)} auseinander. Es ist der Zug, der
nach der \VUgras{Taufe} des Notarssohns aus der Kirche kommt.

%%K t03-c1-12-1 p
12. --- Der kleine Jakob wacht in seiner \VUgras{Wiege} \textsuperscript{(26)} vom
fröhlichen Geläut der Glocken auf, und seine Schwester Magdalene, die den Zug
auch sehen will, bittet ihre Mutter, sie auf der Türschwelle zu kämmen und
anzuziehen.

%%K t03-c1-13-1 p
13. --- Die Mutter ist einverstanden. Sie setzt ihr \VUgras{Kopftuch}
\textsuperscript{(28)} auf, zieht ihr \VUgras{Mieder} \textsuperscript{(29)} und ihre
\VUgras{Schürze} \textsuperscript{(30)} an und setzt sich auf einen kleinen Stuhl vor
die Tür. Magdalene hat nur ihren \VUgras{Unterrock} \textsuperscript{(31)} und ihr
\VUgras{Mieder} \textsuperscript{(32)} an.

%%K t03-c1-14-1 p
14. --- Wie glücklich sie ist! Sie vergisst ihre Lieblingsblumen, ihre
\VUgras{Nelken} \textsuperscript{(34)}, auf denen sich die \VUgras{Schmetterlinge}
\textsuperscript{(35)} niederlassen, ihre \VUgras{Tulpen} \textsuperscript{(36)}, ihre
\VUgras{Maiglöckchen} \textsuperscript{(37)} in den \VUgras{Töpfen}
\textsuperscript{(38)} auf der \VUgras{Mauer} \textsuperscript{(33)} und ihre
\VUgras{Rosen} \textsuperscript{(39)}, die eine so schöne Hecke bilden. Sie schaut
nur auf die schönen Kleider der Damen im Zug.

%%K t03-c1-15-1 p
15. --- Die Dorfjungen achten wenig auf die Kleider. Sie stürzen vor und drängeln
sich, um \VUgras{Bonbons} \textsuperscript{(55)} oder \VUgras{Münzen}
\textsuperscript{(52)} zu bekommen. Einer von ihnen weint \textsuperscript{(40)}.

%%K t03-c1-16-1 p
16. --- Ein anderer ist dem \VUgras{Hund} \textsuperscript{(42)} auf die Pfote
getreten; der läuft jaulend davon und verscheucht die \VUgras{Henne}
\textsuperscript{(43)} mit ihren \VUgras{Küken} \textsuperscript{(44)}.

%%K t03-tit-6 sub
\VUsaut{5.0mm}
\VUpk{10.2pt}{40}{Der Taufzug.}
\VUsaut{4.0mm}

%%K t03-c1-17-1 p
17. --- An der Spitze des Zuges geht die \VUgras{Amme} \textsuperscript{(45)}. Sie
trägt eine schöne \VUgras{Haube} \textsuperscript{(46)} mit langen Bändern. Sie
trägt das \VUgras{Neugeborene} \textsuperscript{(47)}, ganz in \VUgras{Spitzen}
\textsuperscript{(48)} gekleidet.

%%K t03-c1-18-1 p
18. --- Hinter der Amme kommen der \VUgras{Pate} \textsuperscript{(49)} und die
\VUgras{Patin} \textsuperscript{(53)}. Sie verteilen die Bonbons und die Münzen. Der
Pate hat \VUgras{Handschuhe} \textsuperscript{(51)} und einen \VUgras{Zylinder}
\textsuperscript{(50)}. Die Patin hält einen schönen \VUgras{Strauß}
\textsuperscript{(54)} in der Hand.

%%K t03-c1-19-1 p
19. --- Dahinter sehen wir den \VUgras{Onkel} \textsuperscript{(56)} und die
\VUgras{Tante} \textsuperscript{(58)} des Kindes. Die Tante hat einen eleganten
\VUgras{Hut} \textsuperscript{(59)}, der Onkel einen schwarzen \VUgras{Gehrock}
\textsuperscript{(57)} und eine weiße Weste. Dann kommen der \VUgras{Vetter}
\textsuperscript{(60)} und seine \VUgras{Base} \textsuperscript{(61)}. Sie hält die
\VUgras{Schwester} \textsuperscript{(62)} des Neugeborenen an der Hand, die kleine
Bertha.

%%K t03-c1-20-1 p
20. --- Berthas anderer \VUgras{Bruder} \textsuperscript{(63)} geht dahinter. Er gibt
einem \VUgras{Verwandten} \textsuperscript{(64)} die Hand. Die \VUgras{Freunde}
\textsuperscript{(65)} der Familie kommen zuletzt.

%%K t03-tit-7 sub
\VUsaut{4.6mm}
\VUpk{10.2pt}{40}{Die Zuschauer.}
\VUsaut{3.4mm}

%%K t03-c1-21-1 p
21. --- Neben dem Zug stützt sich ein \VUgras{Bettler} \textsuperscript{(66)} auf
seine \VUgras{Krücke} \textsuperscript{(67)}. Er bittet um Geld.

%%K t03-c1-22-1 p
22. --- Auf der anderen Seite steht ein sehr \VUgras{höflicher}
\textsuperscript{(68)} kleiner Junge. Er zieht die Mütze und wünscht den Leuten,
die er kennt, einen \VUgras{„guten Morgen“}.

%%K t03-c1-23-1 p
23. --- Die Dorfbewohner sind aus ihren Häusern gekommen, um den Taufzug
vorbeiziehen zu sehen. Wir sehen einen \VUgras{Bauern} \textsuperscript{(69)} mit
seiner \VUgras{Zipfelmütze} und neben ihm eine \VUgras{Bäuerin}
\textsuperscript{(70)}. Dann eine \VUgras{alte Frau} \textsuperscript{(71)}, die sich
auf ihren \VUgras{Stock} \textsuperscript{(72)} stützt. Zuletzt den \VUgras{Müller}
\textsuperscript{(73)} mit seinem breiten \VUgras{Filzhut} \textsuperscript{(74)} und
eine junge Bäuerin in \VUgras{Holzschuhen} \textsuperscript{(75)}, die ihrem Kind
das Laufen beibringt.

%%K t03-c1-24-1 p
24. --- Es ist eine sehr lebhafte Szene.
\VUsaut{4.0mm}

%%K t03-tit-8 sub
\VUcentre{12.0pt}{0}{\textit{Zweite Szene.}}
\VUsaut{3.0mm}

%%K t03-tit-9 sub
\VUpk{10.2pt}{40}{Ein Feiertag.}
\VUsaut{2.4mm}

%%K t03-tit-10 sub
\VUpk{10.2pt}{40}{Wassersport und Regatta.}
\VUsaut{3.0mm}

%%K t03-c2-01-1 p
1. --- Der \VUgras{Sommer} ist da. Die brennende Sonne des \VUgras{Juni}
(Juli oder \VUgras{August}) lockt die jungen Leute zum Baden und zum Rudern.

%%K t03-c2-02-1 p
2. --- Am rechten Ufer des Flusses ziehen sich die Ruderer aus, um am Wassersport
und an der Regatta teilzunehmen.

%%K t03-c2-03-1 p
3. --- Einer von ihnen zieht seine \VUgras{Schuhe} \textsuperscript{(1)} aus; er
schnürt (knöpft) sie auf. Seine beiden Freunde sind schon fertig. Sie haben nur
noch ihr \VUgras{Trikot} \textsuperscript{(2)} und ihre \VUgras{Badehose}
\textsuperscript{(3)} an. Sie plaudern, während sie auf ihre Freunde warten. Sie
reden vom Jahrmarkt und vom Fest am anderen Ufer.

%%K t03-c2-04-1 p
4. --- Neben ihnen liegen auf dem Boden die \VUgras{Kleider} ihrer Freunde: ein
\VUgras{Paar} \VUgras{Stiefel} \textsuperscript{(4)}, \VUgras{Schuhe}
\textsuperscript{(5)}, \VUgras{Socken} \textsuperscript{(6)}, \VUgras{Filz-}
\textsuperscript{(7)} und \VUgras{Strohhüte} \textsuperscript{(8)} und eine
\VUgras{Krawatte} \textsuperscript{(9)}.

%%K t03-c2-05-1 p
5. --- Einer der jungen Männer hat seine \VUgras{Jacke} \textsuperscript{(10)}
sorgfältig zusammengelegt. Er legt sie auf ein Handtuch, in das sein Nachbar
gleich beide \VUgras{Anzüge} einwickeln wird.

%%K t03-c2-06-1 p
6. --- Ein sechster Junge ist spät dran. Er hat noch seine \VUgras{Mütze}
\textsuperscript{(11)} auf dem Kopf. Er hat weder sein \VUgras{Hemd}
\textsuperscript{(12)} noch seinen \VUgras{Kragen} \textsuperscript{(13)} noch seine
\VUgras{Manschetten} \textsuperscript{(14)} ausgezogen. Er hält seine
\VUgras{Weste} \textsuperscript{(17)} in der Hand und trägt noch seine
\VUgras{Hose} \textsuperscript{(16)} und seine \VUgras{Hosenträger}
\textsuperscript{(15)}. Zum nächsten Rennen ist er sicher nicht bereit.

%%K t03-c2-07-1 p
7. --- Neben den jungen Männern führt ein Weg zwischen \VUgras{Disteln}
\textsuperscript{(18)} zum Fluss hinunter, wo der alte Jakob im \VUgras{Schilf}
\textsuperscript{(19)} das \VUgras{Feuerwerk} \textsuperscript{(20)} vorbereitet, das
am Abend abgebrannt wird.

%%K t03-tit-11 sub
\VUsaut{4.4mm}
\VUpk{10.2pt}{40}{Das Rennen.}
\VUsaut{3.4mm}

%%K t03-c2-08-1 p
8. --- Auf dem Fluss sitzen ein paar Damen unter ihren Sonnenschirmen in einem
\VUgras{Ruderboot} \textsuperscript{(21)}. Sie verfolgen das Rennen, das sehr
spannend ist. In jedem \VUgras{Rennboot} \textsuperscript{(22)} ziehen vier
\VUgras{Ruderer} \textsuperscript{(24)} mit aller Kraft, während der
\VUgras{Steuermann} \textsuperscript{(26)} das \VUgras{Steuer}
\textsuperscript{(25)} hält und das leichte Boot lenkt. Die \VUgras{Riemen}
\textsuperscript{(23)} schlagen im Takt ins Wasser, und die leichten Boote fliegen
mit schwindelerregender Geschwindigkeit dahin.

%%K t03-c2-09-1 p
9. --- Ein paar junge Männer wetteifern neben dem Brückenpfeiler um den Preis des
\VUgras{Klettermasts} \textsuperscript{(28)}. Einer von ihnen schiebt sich vorsichtig
auf der federnden, glatten Stange vorwärts, um an das \VUgras{Fähnchen}
\textsuperscript{(29)} am Ende zu kommen. Sein glückloser \VUgras{Gegner}
\textsuperscript{(30)} ist ins Wasser gefallen. Er schwimmt ans Ufer zurück.

%%K t03-c2-10-1 p
10. --- Ältere Burschen treiben \VUgras{Wasserstechen} \textsuperscript{(32)} neben
dem \VUgras{Kai} \textsuperscript{(33)}. Eine große \VUgras{Menge}
\textsuperscript{(31)} schaut all diesen Spielen zu, an die \VUgras{Brüstung} der
\VUgras{Brücke} \textsuperscript{(27)} gelehnt und dicht am \VUgras{Ufer} gedrängt.
Ein paar Zuschauer drehen sich allerdings um, um dem Spiel an der
\VUgras{Kletterstange} \textsuperscript{(45)} zu folgen oder den
\VUgras{Zug des Bürgermeisters} auf dem Weg zur \VUgras{Preisverleihung} zu
bewundern.

%%K t03-c2-11-1 p
11. --- Zuerst kommen ein paar \VUgras{Gassenjungen} \textsuperscript{(35)}, die der
\VUgras{Kapelle} \textsuperscript{(36)} vorauslaufen; dann der
\VUgras{Bürgermeister} \textsuperscript{(37)}, seine Beigeordneten und der
\VUgras{Gemeinderat} des Städtchens. Ganz zuletzt marschieren die
\VUgras{Feuerwehrleute} \textsuperscript{(38)}.

%%K t03-tit-12 sub
\VUsaut{5.0mm}
\VUpk{10.2pt}{40}{Der Jahrmarkt.}
\VUsaut{3.6mm}

%%K t03-c2-12-1 p
12. --- Auf dem \VUgras{Festplatz} \textsuperscript{(39)} drängt sich eine große
Menge. Man kommt, um die verschiedenen \VUgras{Buden} zu sehen: das
\VUgras{Karussell} \textsuperscript{(40)}, den \VUgras{Zirkus} \textsuperscript{(41)},
wo die Vorstellung schon begonnen hat, und den \VUgras{Tanzboden}
\textsuperscript{(42)}, wo die jungen Leute der Stadt sich am \VUgras{Tanzen}
freuen.

%%K t03-c2-13-1 p
13. --- Ganz hinten sehen wir die \VUgras{Arena} \textsuperscript{(44)}, wo der
tapfere spanische \VUgras{Matador} gegen den wütenden \VUgras{Stier}
\textsuperscript{(45)} kämpft; dann die \VUgras{Achterbahn} \textsuperscript{(49)} und
die \VUgras{Schießbude} \textsuperscript{(46)}, wo man den Umgang mit
\VUgras{Schusswaffen} und das Schießen auf die \VUgras{Scheibe} übt.

%%K t03-c2-14-1 p
14. --- Gleich daneben steht das \VUgras{Jahrmarktstheater} \textsuperscript{(47)}, wo
\VUgras{Komödie} und \VUgras{Drama} gespielt werden. Ganz in der Nähe steht die
Bude des \VUgras{Riesen} \textsuperscript{(48)} und des \VUgras{Zwerges}. Die
\VUgras{Größe} des einen beträgt zwei Meter fünfzehn; der andere ist kaum so groß
wie ein Kind.

%%K t03-c2-15-1 p
15. --- Zuletzt kommt die große \VUgras{Menagerie} \textsuperscript{(50)} mit ihren
\VUgras{wilden Tieren}, ihrem \VUgras{Tiger} \textsuperscript{(51)} mit den
kräftigen \VUgras{Krallen}, ihrem \VUgras{Löwen} \textsuperscript{(52)} mit der
dichten \VUgras{Mähne}, ihren beiden \VUgras{Giraffen} \textsuperscript{(53)} und
ihrem \VUgras{Elefanten} \textsuperscript{(54)}, der seinen \VUgras{Rüssel} so
geschickt gebraucht.

%%K t03-c2-16-1 p
16. --- Der \VUgras{Bändiger} \textsuperscript{(55)}, der diese Tiere abrichtet, steht
an der Tür der Bude.
\VUsaut{6.0mm}
\VUfilet{22mm}
